% ============================================================
% Model: 空穴概念与半导体输运
% 合并自 hole-concept.tex（一维六要素定义）与 gaas-hole-drude.tex（三维半导体输运）
% ============================================================

\section{空穴概念与半导体输运}
\label{sec:hole-concept-and-transport}

\begin{tipbox}[关键词标签]
空穴 · 有效质量 · 准动量 · 满带 · 价带顶 · Drude模型 · 电导率 · p型半导体 · GaAs
\end{tipbox}

% --------------------------------------------------------
% 1. 模型定义框
% --------------------------------------------------------
\begin{modelbox}[模型：从空穴定义到半导体输运]
\textbf{物理情境}：当能带被电子几乎填满时，少数电子从价带顶被激发到导带，留下的空态可等效为带正电的准粒子——\key{空穴}。在电场作用下，空穴像经典带正电粒子一样漂移，可用 Drude 模型计算其输运性质。

\textbf{空穴的物理定义}：
\begin{itemize}[nosep]
    \item 空穴是价带中大量电子集体运动的等效描述
    \item 空穴波矢 $\mathbf{k}_h=-\mathbf{k}_e$（相对于价带顶）
    \item 空穴能量 $\varepsilon_h(\mathbf{k})=-\varepsilon_e(\mathbf{k})$（以价带顶为能量零点）
    \item 空穴有效质量 $m_h^*=-m_e^*>0$（价带电子有效质量为负）
    \item 空穴速度 $\mathbf{v}_h=\nabla_{\mathbf{k}}\varepsilon_h/\hbar=\nabla_{\mathbf{k}}\varepsilon_e/\hbar=\mathbf{v}_e$
\end{itemize}

\textbf{核心对应关系（六要素）}：
\begin{center}
\begin{tabular}{ccc}
\toprule
\textbf{物理量} & \textbf{电子} & \textbf{空穴} \\
\midrule
波矢 & $k_e$ & $k_h=-k_e$ \\
电荷 & $-e$ & $+e$ \\
能量 & $E_e(k_e)$ & $E_h=-E_e(k_e)$ \\
速度 & $v_e=\frac{1}{\hbar}\frac{\dd E_e}{\dd k_e}$ & $v_h=v_e$ \\
有效质量 & $m_e^*=\hbar^2\bigl(\frac{\dd^2E_e}{\dd k_e^2}\bigr)^{-1}$ & $m_h^*=-m_e^*$ \\
准动量 & $p_e=\hbar k_e$ & $p_h=\hbar k_h=-p_e$ \\
\bottomrule
\end{tabular}
\end{center}

\textbf{GaAs 价带特征}：
\begin{itemize}[nosep]
    \item 价带顶在 $\Gamma$ 点（$\mathbf{k}=0$），三重简并（考虑自旋为六重）
    \item 轻空穴带（LH）：$m_{lh}^*\approx0.08m_0$
    \item 重空穴带（HH）：$m_{hh}^*\approx0.5m_0$
    \item 自旋--轨道耦合分裂带（SO）：$\Delta_{SO}\approx0.34\,$eV
\end{itemize}
\end{modelbox}

% --------------------------------------------------------
% 2. 关键公式框
% --------------------------------------------------------
\begin{formulabox}[核心公式]
\begin{enumerate}[label=\textbf{(\arabic*)}]
    \item \textbf{空穴波矢}：
    \[
    \mathbf{k}_h=-\mathbf{k}_e
    \]
    \texteq{满带总动量为零，移走电子后系统动量等于 $-\mathbf{k}_e$}

    \item \textbf{空穴有效质量}：
    \[
    m_h^*=-m_e^*=-\hbar^2\Bigl(\frac{\partial^2E_e}{\partial k_e^2}\Bigr)^{-1}
    \]
    \texteq{能带顶 $m_e^*<0$，故 $m_h^*>0$，空穴表现如同正质量正电荷粒子}

    \item \textbf{空穴速度}：
    \[
    \mathbf{v}_h=\mathbf{v}_e=\frac{1}{\hbar}\nabla_{\mathbf{k}}E_e\big|_{\mathbf{k}_e}
    \]
    \texteq{速度不变，因为 $E_h(\mathbf{k}_h)=-E_e(-\mathbf{k}_h)$ 对 $\mathbf{k}_h$ 求导后符号抵消}

    \item \textbf{空穴运动方程}：
    \[
    \hbar\frac{d\mathbf{k}_h}{dt}=+e(\mathbf{E}+\mathbf{v}_h\times\mathbf{B})
    \]
    \texteq{注意正号！空穴带电荷 $+e$，受力方向与电场相同}

    \item \textbf{空穴漂移速度}：
    \[
    \mathbf{v}_d=\frac{e\tau}{m_h^*}\mathbf{E}=\mu_h\mathbf{E}
    \]
    \texteq{$\mu_h=e\tau/m_h^*$ 为空穴迁移率}

    \item \textbf{Drude 电导率（p 型）}：
    \[
    \sigma_p=\frac{n_h e^2\tau}{m_h^*}=n_h e\mu_h
    \]
    \texteq{$n_h$ 为空穴浓度，$m_h^*$ 为空穴有效质量}

    \item \textbf{重/轻空穴混合输运}：
    \[
    \sigma_p=e^2\tau\Bigl(\frac{n_{hh}}{m_{hh}^*}+\frac{n_{lh}}{m_{lh}^*}\Bigr)
    \]
    \texteq{总电导率为各子带贡献之和，重空穴通常占主导（$m_{hh}^*$ 大但态密度高）}

    \item \textbf{Hall 系数（p 型）}：
    \[
    R_H=\frac{1}{n_h e}
    \]
    \texteq{p 型半导体 $R_H>0$，这是判定导电类型的金标准}
\end{enumerate}
\end{formulabox}

% --------------------------------------------------------
% 3. 训练点框
% --------------------------------------------------------
\begin{drillbox}[训练点与考法变体]
\trainingpoint{1} \textbf{六要素全算}：给定 $E(k)$ 和移走电子的位置 $k_e$，依次计算 $k_h, m_h^*, v_h, p_h, E_h, \dd p_h/\dd t$。

\trainingpoint{2} \textbf{有效质量符号判断}：给定 $E(k)$ 曲线，判断某 $k$ 点电子有效质量正负，进而确定空穴有效质量正负。只有 $m_h^*>0$ 的空穴才是稳定的准粒子。

\trainingpoint{3} \textbf{速度方向判断}：已知电子速度方向，判断空穴速度方向（相同），但电流方向由正电荷运动决定。

\trainingpoint{4} \textbf{空穴运动方程的符号}：推导空穴在电场和磁场中的运动方程，特别注意空穴电荷为 $+e$，洛伦兹力公式中的符号与电子相反。

\trainingpoint{5} \textbf{p 型与 n 型对比}：给定相同掺杂浓度和散射时间，比较 n 型和 p 型 GaAs 的电导率和迁移率（$\mu_n\approx8500\,\text{cm}^2/\text{Vs}$, $\mu_p\approx400\,\text{cm}^2/\text{Vs}$）。

\trainingpoint{6} \textbf{重/轻空穴浓度比}：在热平衡下，$n_{hh}/n_{lh}\propto(m_{hh}^*/m_{lh}^*)^{3/2}$，重空穴浓度占优，验证总电导率由重空穴主导。

\trainingpoint{7} \textbf{外场漂移与霍尔效应}：在电场/磁场中，分别写出电子和空穴的运动方程，验证霍尔效应中 n 型与 p 型的霍尔系数符号相反。
\end{drillbox}

% --------------------------------------------------------
% 4. 解题技巧框
% --------------------------------------------------------
\begin{tipbox}[快速识别与解题技巧]
\begin{itemize}
    \item \examnote{识别标志}：题干出现``满带''''移走电子''''产生空穴''''价带顶''''p 型''''Drude 模型''，立即定位本模型。
    \item \examnote{六要素速查卡}：
    \begin{center}
    \fbox{\parbox{0.9\linewidth}{\centering
    \textbf{空穴 = ``电子的反面''，但速度保留}\\[2pt]
    $k_h=-k_e$\quad$E_h=-E_e$\quad$p_h=-p_e$\\[2pt]
    $m_h^*=-m_e^*$
    \quad$v_h=v_e$
    \quad$\dot{p}_h=+e\mathcal{E}$
    }}
    \end{center}
    \item \examnote{空穴 vs 电子速记}：
    \begin{center}
    \begin{tabular}{lll}
    \toprule
    \textbf{物理量} & \textbf{电子} & \textbf{空穴} \\
    \midrule
    电荷 & $-e$ & $+e$ \\
    波矢（价带顶） & $\mathbf{k}_e$ & $\mathbf{k}_h=-\mathbf{k}_e$ \\
    能量 & $\varepsilon_e$ & $\varepsilon_h=-\varepsilon_e$ \\
    有效质量 & $m_e^*$（可负） & $m_h^*=-m_e^*>0$ \\
    电场受力 & $-e\mathbf{E}$ & $+e\mathbf{E}$ \\
    漂移速度 & $-e\tau\mathbf{E}/m_e^*$ & $+e\tau\mathbf{E}/m_h^*$ \\
    Hall 系数 & $-1/ne$ & $+1/n_h e$ \\
    \bottomrule
    \end{tabular}
    \end{center}
    \item \examnote{有效质量记忆}：能带底 $m_e^*>0$（电子），$m_h^*<0$（无物理空穴）；能带顶 $m_e^*<0$（电子），$m_h^*>0$（真实空穴）。
    \item \examnote{Drude 模型三大假设}：
    \begin{enumerate}[nosep]
        \item 独立粒子近似（电子/空穴间无相互作用）
        \item 自由加速 + 瞬时碰撞（碰撞后速度随机化）
        \item 弛豫时间近似（平均自由程 $\ell=v_F\tau$）
    \end{enumerate}
    \item \examnote{GaAs 参数速查}：
    \begin{itemize}[nosep]
        \item $m_{hh}^*=0.50m_0$, $m_{lh}^*=0.08m_0$
        \item $\varepsilon_g=1.42\,$eV（300K，直接带隙）
        \item $\mu_n=8500\,\text{cm}^2/\text{Vs}$, $\mu_p=400\,\text{cm}^2/\text{Vs}$
    \end{itemize}
\end{itemize}
\end{tipbox}

% --------------------------------------------------------
% 5. 易错点框
% --------------------------------------------------------
\begin{errorbox}{常见失误与陷阱}
\begin{itemize}
    \item \textbf{速度方向混淆}：电子从价带顶被激发后向更高 $k$ 运动（若 $k_e>0$），空穴 $k_h=-k_e<0$。但 $v_h=v_e$，两者速度方向相同。电流方向由正电荷运动决定，故空穴电流与电子电流同向。
    \item \textbf{空穴有效质量符号}：价带电子有效质量 $m_e^*<0$（曲率为负），但空穴有效质量 $m_h^*=-m_e^*>0$。在电导率公式中直接用 $m_h^*$ 即可，无需额外负号。
    \item \textbf{能量符号}：$E_h=-E_e$ 是以满带为零点的结果。若 $E_e<0$（价带顶通常取为能量零点以下），则 $E_h>0$，表示移走电子需要输入能量。
    \item \textbf{有效质量推导}：$m_h^*=-m_e^*$ 的严格推导需要用到 $k_h=-k_e$ 和链式法则：
    \[
    \frac{\dd^2E_h}{\dd k_h^2}=\frac{\dd^2(-E_e)}{\dd(-k_e)^2}=\frac{\dd^2(-E_e)}{\dd k_e^2}=-\frac{\dd^2E_e}{\dd k_e^2}.
    \]
    不要凭直觉直接写正负号。
    \item \textbf{准动量变化率}：电子 $\dot{p}_e=-e\mathcal{E}$，空穴 $\dot{p}_h=+e\mathcal{E}$。注意电场力对正/负电荷方向相反。
    \item \textbf{空穴速度}：$\mathbf{v}_h=\mathbf{v}_e$！虽然波矢相反，但能量也相反，所以群速度相同。不要误认为空穴速度也与电子相反。
    \item \textbf{重空穴主导电导}：虽然重空穴迁移率低（$\mu_{hh}<\mu_{lh}$），但重空穴浓度高（$n_{hh}\gg n_{lh}$），总电导率通常由重空穴贡献主导。
    \item \textbf{一维 vs 三维}：一维中 $m^*$ 是标量；三维中 $m^*$ 是张量，$1/m^*_{\alpha\beta}=\frac{1}{\hbar^2}\frac{\partial^2E}{\partial k_\alpha\partial k_\beta}$。
\end{itemize}
\end{errorbox}

% --------------------------------------------------------
% 6. 物理图像与关联模型
% --------------------------------------------------------
\begin{insightbox}[物理图像与模型关联]
\textbf{物理图像}：

把价带想象成一个坐满人的剧场（满带）：
\begin{itemize}[nosep]
    \item \textbf{电子}：一个观众从座位 $k_e$ 站起来走到后台（被激发到导带）。
    \item \textbf{空穴}：剩下的观众看到的不是``某人走了''，而是``某个座位上出现了一个空位''。这个空位的运动（被后方观众填补）比追踪每个观众的位置更容易描述。
    \item \textbf{正电荷}：空位向前移动时，实际上是后方的人依次向前补位——等效于一个``正电荷''在向前移动。
    \item \textbf{速度}：空位移动的速度与观众实际移动的速度相同——$v_h=v_e$。
\end{itemize}

\textbf{与相似模型的对比}：
\begin{center}
\begin{tabular}{lllll}
\toprule
\textbf{对比维度} & \textbf{导带电子} & \textbf{价带空穴} & \textbf{n 型半导体} & \textbf{p 型半导体} \\
\midrule
电荷 & $-e$ & $+e$ & $-e$ & $+e$ \\
能带位置 & 导带底 & 价带顶 & 导带底 & 价带顶 \\
有效质量 & $m_e^*>0$ & $m_h^*>0$ & $m_e^*\sim0.01$--$0.1m_0$ & $m_h^*\sim0.1$--$1m_0$ \\
运动方程 & $\dot{p}=-e\mathcal{E}$ & $\dot{p}=+e\mathcal{E}$ & $\dot{p}=-e\mathcal{E}$ & $\dot{p}=+e\mathcal{E}$ \\
Hall 系数 & $-1/ne$ & $+1/n_h e$ & $<0$ & $>0$ \\
\bottomrule
\end{tabular}
\end{center}

\textbf{推广方向}：
\begin{itemize}[nosep]
    \item 磁场中的空穴：朗道能级 $E_n=\hbar\omega_c(n+1/2)$，$\omega_c=eB/m_h^*$
    \item 激子：电子--空穴束缚态，类似氢原子，能量 $E_n=-\frac{\mu e^4}{2\hbar^2\varepsilon^2n^2}$
    \item 极化子：电子/空穴与晶格极化场耦合形成的准粒子
    \item 量子阱中的空穴：量子限制使重/轻空穴带劈裂，改变有效质量和态密度
    \item 自旋--轨道耦合效应：Dresselhaus/Rashba 自旋劈裂影响空穴自旋输运
\end{itemize}
\end{insightbox}

% --------------------------------------------------------
% 7. 推导附录
% --------------------------------------------------------
\begin{derivationbox}[推导细节（自我检验用）]
\textbf{Step 1：空穴速度}

由 $k_h=-k_e$，$E_h(k_h)=-E_e(-k_h)=-E_e(k_e)$。

\[
v_h=\frac{1}{\hbar}\frac{\dd E_h}{\dd k_h}=\frac{1}{\hbar}\frac{\dd(-E_e)}{\dd(-k_e)}=\frac{1}{\hbar}\frac{\dd E_e}{\dd k_e}=v_e.
\]

\textbf{Step 2：空穴有效质量}

\[
\frac{1}{m_h^*}=\frac{1}{\hbar^2}\frac{\dd^2E_h}{\dd k_h^2}
=\frac{1}{\hbar^2}\frac{\dd^2(-E_e)}{\dd(-k_e)^2}
=-\frac{1}{\hbar^2}\frac{\dd^2E_e}{\dd k_e^2}
=-\frac{1}{m_e^*}.
\]

故 $m_h^*=-m_e^*$。

\textbf{Step 3：外场下运动}

电子半经典方程：$\hbar\dot{k}_e=-e\mathcal{E}$。

空穴：$\hbar\dot{k}_h=\hbar\frac{\dd}{\dd t}(-k_e)=-\hbar\dot{k}_e=+e\mathcal{E}$。

即 $\dot{p}_h=+e\mathcal{E}$。

\textbf{Step 4：Drude 电导率推导}

弛豫时间近似：
\[
\frac{d\mathbf{v}_d}{dt}=\frac{e\mathbf{E}}{m_h^*}-\frac{\mathbf{v}_d}{\tau}=0\quad\text{(稳态)}
\]
\[
\Rightarrow\quad\mathbf{v}_d=\frac{e\tau}{m_h^*}\mathbf{E}=\mu_h\mathbf{E}.
\]

电流密度：
\[
\mathbf{j}=n_h(+e)\mathbf{v}_d=\frac{n_h e^2\tau}{m_h^*}\mathbf{E}=\sigma_p\mathbf{E}.
\]

\textbf{Step 5：重/轻空穴混合}

GaAs 价带顶附近两支空穴带，总电流：
\[
\mathbf{j}=\mathbf{j}_{hh}+\mathbf{j}_{lh}=e^2\tau\Bigl(\frac{n_{hh}}{m_{hh}^*}+\frac{n_{lh}}{m_{lh}^*}\Bigr)\mathbf{E}.
\]

浓度比由有效态密度决定：
\[
\frac{n_{hh}}{n_{lh}}=\Bigl(\frac{m_{hh}^*}{m_{lh}^*}\Bigr)^{3/2}\Bigl(\frac{e^{-\varepsilon_{hh}/k_BT}}{e^{-\varepsilon_{lh}/k_BT}}\Bigr).
\]

在带边（两支能量简并），$n_{hh}/n_{lh}\approx(0.5/0.08)^{3/2}\approx16$，重空穴占主导。

\textbf{Step 6：一维具体算例}

设 $E(k)=\frac{\hbar^2k_i^2}{6m}-\frac{3\hbar^2k^2}{m}$，$k_i=\frac{1}{8a}$，$a=0.314\,\text{nm}$。

移走 $k_e=k_i$ 处电子：
\begin{itemize}[nosep]
    \item $k_h=-k_i=-3.98\times10^8\,\text{m}^{-1}$
    \item $\frac{\dd^2E}{\dd k^2}=-\frac{6\hbar^2}{m}$，$m_e^*=-\frac{m}{24\pi^2}$，$m_h^*=\frac{m}{24\pi^2}$
    \item $v_e=\frac{1}{\hbar}\frac{\dd E}{\dd k}\big|_{k_i}=-\frac{12\pi\hbar k_i}{m}=v_h$
    \item $p_h=-\hbar k_i=-\frac{h}{16\pi a}$
    \item $E_h=-E(k_i)=\frac{17\hbar^2k_i^2}{6m}$
    \item $\frac{\dd p_h}{\dd t}=+e\mathcal{E}$
\end{itemize}
\end{derivationbox}

\vspace{1em}
